\documentclass{report}
\usepackage{amsfonts, amsmath}

\newcommand\N{{\mathbb N}}
\newcommand\Q{{\mathbb Q}}
\newcommand\R{{\mathbb R}}
\newcommand\Z{{\mathbb Z}}

\pagestyle{empty}
\begin{document}
\large
\centerline{\Huge \bf Analisi Matematica I}
\vskip.2cm
\centerline{\Huge  Quarto Appello (09-09-2002)}
\renewcommand{\thefootnote}{ } 
\footnote{\sl Avete 3 ore di tempo. Ogni esercizio
vale sei punti. La sufficienza si ottiene con un punteggio $\geq$
18. {\bf Solo le risposte chiaramente giustificate saranno prese in
considerazione.} {\scshape Elaborati illeggibili o confusi verranno ignorati.}
{\sffamily La
sintesi sar\`a particolarmente apprezzata.}}
\renewcommand{\thefootnote}{\arabic{footnote}} 

\vskip1cm
\begin{enumerate}
\item Dire quale dei seguenti numeri \`e il maggiore
\[
\log_3 9\;;\quad 2\;;\quad \log_2 5.
\]
\item Si trovino i divisori comuni di $q=1757$ e $p=4819$.

\item Si consideri la funzione
\[
f(x):=\frac 12 x+\frac 14\sin x.
\]
Si verifichi che $f([0,1])\subset [0,1]$ e che
$\sup\limits_{x\in[0,1]}|f'(x)|\leq \frac 34$.\newline Usando tali fatti si dimostri
che
\[
\lim_{n\to\infty}f^n(0.5)=0.
\]
\item Si consideri la funzione
\[
f(x):=\ln(|x-1|+x).
\]
Se ne determini il dominio e si disegni il grafico. Si trovi il numero
delle soluzioni dell'equazione $f(x)=3$.
\item Si studi la convergenza della serie
\[
\sum_{n=2}^\infty \ln\left(\sqrt{\frac{n^3+n}{n^3-1}}\right).
\]
\item Si consideri una circonferenza di raggio 1. Tra tutti i triangoli
inscritti e con un lato di lunghezza 2, si caratterizzino quelli di area
massima.
\end{enumerate}
\end{document}
